\chapter{Food Groups and a Balanced Diet}

\section*{Definition and Overview}
A balanced diet is one that provides the right types and amounts of food to meet the body's needs for energy, protein, vitamins, minerals, and fibre. The Australian Dietary Guidelines organise foods into five core groups: vegetables and legumes; fruit; grain foods, mostly wholegrain; lean meats, poultry, fish, eggs, tofu, nuts, and seeds; and dairy products or alternatives. Eating a variety of foods from each group, in the recommended amounts, is the foundation of good health and the prevention of chronic disease. This chapter explains the food groups, how much of each is recommended, and how to build a balanced plate for every meal.

\section*{Epidemiology}
Most Australians do not eat a fully balanced diet. National surveys show that fewer than one in ten adults eats the recommended serves of vegetables, and many eat too little fruit, dairy, and wholegrains while consuming too many discretionary foods such as takeaway, confectionery, and sugary drinks. Poor dietary patterns contribute to the high rates of obesity, heart disease, type 2 diabetes, and some cancers in Australia, and dietary risk factors are among the leading causes of preventable death. Improving alignment with the food group recommendations is a major public health priority, and even small changes produce measurable benefit.

\section*{Aetiology and Pathophysiology}
Each food group makes a distinct contribution to health.
\begin{itemize}
    \item \textbf{Vegetables and legumes:} rich in vitamins, minerals, fibre, and protective plant compounds; associated with lower rates of heart disease, stroke, and some cancers; recommended at least 5 serves a day
    \item \textbf{Fruit:} provides vitamin C, fibre, and antioxidants; recommended 2 serves a day
    \item \textbf{Grain foods:} wholegrains provide carbohydrate for energy, B vitamins, and fibre, and are associated with lower risk of heart disease and diabetes; choose mostly wholegrain
    \item \textbf{Lean meats, poultry, fish, eggs, tofu, nuts, and seeds:} provide protein, iron, zinc, and omega-3 fats in oily fish; some red meat, in moderation
    \item \textbf{Dairy and alternatives:} provide calcium, protein, and iodine; important for bone health; reduced-fat versions are recommended
    \item \textbf{Discretionary foods:} foods high in saturated fat, added sugar, and salt, such as cakes, chips, and sugary drinks, provide energy with little nutrition and should be limited
    \item \textbf{Balance and variety:} no single food provides everything, so variety across groups meets all nutrient needs
    \item \textbf{Portions:} serve sizes within each group are designed so that total intake meets energy needs
    \item \textbf{Fluids:} water is the preferred drink; sugary drinks contribute excess energy
\end{itemize}

\section*{Clinical Features}
A balanced diet has recognisable features.
\begin{itemize}
    \item Vegetables at most meals, making up about half the plate
    \item Fruit as part of breakfast and snacks
    \item Wholegrain bread, rice, pasta, and cereals as the grain choices
    \item A serve of protein food at each main meal
    \item Dairy or a calcium-fortified alternative daily
    \item Water as the main drink, with limited juice and no sugary drinks
    \item Discretionary foods limited to sometimes foods, not everyday foods
    \item Adequate energy for growth and activity, without excess weight gain
\end{itemize}

\section*{Differential Diagnosis}
Dietary advice must be tailored to individuals and conditions.
\begin{itemize}
    \item Pregnancy and breastfeeding, which increase needs for iron, folate, and calcium
    \item Children and adolescents, with growth-related needs
    \item Older adults, who need adequate protein and calcium and benefit from nutrient-dense foods
    \item Diabetes, which requires attention to carbohydrate amounts and types
    \item Heart disease, kidney disease, and food allergies, which require specific modifications
    \item Vegetarian and vegan diets, which need planning to ensure iron, B12, and calcium
    \item The guidelines are a framework, not a rigid prescription, and individual advice comes from a dietitian
\end{itemize}

\section*{When to Seek Medical Care}
Most people can improve their diet without medical care, but a doctor or dietitian should be consulted for dietary needs related to medical conditions, pregnancy, or children with poor growth, and for weight management that is not responding.

\textbf{Call 000 (or local emergency services) for:}
\begin{itemize}
    \item Any medical emergency, including symptoms of severe malnutrition or dehydration
    \item Difficulty swallowing or choking
    \item Sudden severe symptoms in someone with a dietary-related condition
\end{itemize}

\section*{Diagnosis and Investigations}
\begin{itemize}
    \item \textbf{Dietary assessment:} a dietitian reviews usual intake against the food group targets
    \item \textbf{Growth monitoring:} weight and height plotted over time in children
    \item \textbf{Blood tests:} for specific nutrients, such as iron, B12, vitamin D, or calcium, when deficiency is suspected
    \item \textbf{Screening:} blood pressure, blood sugar, and cholesterol for people with dietary risk factors
    \item \textbf{Individualised planning:} advice tailored to age, activity, culture, and medical conditions
\end{itemize}

\section*{Management}
\begin{itemize}
    \item \textbf{Build meals around the groups:} aim for vegetables, a wholegrain, and a protein food at each main meal, with fruit and dairy across the day
    \item \textbf{Use the recommended serves:} 5+ vegetables, 2 fruit, 4--6 grains, 2.5--3 protein, and 2--4 dairy serves daily for adults, adjusted for age and sex
    \item \textbf{Choose wholegrains:} wholegrain bread, brown rice, oats, and quinoa
    \item \textbf{Limit discretionary foods:} reduce takeaway, confectionery, chips, and sugary drinks to occasional treats
    \item \textbf{Plan meals and snacks:} shopping and preparing meals ahead makes healthy choices easier
    \item \textbf{Cook more at home:} home cooking allows control of salt, fat, and sugar
    \item \textbf{Read labels:} compare products for fat, sugar, and salt
    \item \textbf{Seek professional advice:} a dietitian for individualised plans and medical conditions
\end{itemize}

\section*{Complications}
\begin{itemize}
    \item Nutrient deficiencies, including iron, calcium, and iodine, from unbalanced diets
    \item Obesity and its complications from excess discretionary foods
    \item Heart disease, type 2 diabetes, and some cancers from long-term poor eating patterns
    \item Poor growth and development in children
    \item Osteoporosis from inadequate calcium and vitamin D
    \item Fatigue and reduced immunity from insufficient nutrition
\end{itemize}

\section*{Prognosis and Outcomes}
Shifting towards a balanced diet produces measurable improvements in health: better weight control, lower blood pressure and cholesterol, improved blood sugar, and reduced long-term risk of chronic disease. The benefits accrue at any age, and changes made gradually are more likely to last. People who follow the food group guidance consistently have lower rates of heart disease, diabetes, and cancer, and better energy and wellbeing. A balanced diet, combined with regular activity, is the most powerful single investment in long-term health.

\section*{Prevention and Self-Care}
\begin{itemize}
    \item Aim for the recommended serves from each food group every day
    \item Fill half the plate with vegetables at main meals
    \item Choose wholegrain options over refined
    \item Drink water and limit sugary drinks
    \item Limit takeaway, snacks, and confectionery
    \item Plan meals and shop with a list
    \item Involve children in choosing and preparing food
    \item See a dietitian for individualised advice, especially with medical conditions
\end{itemize}

\section*{Sources and Further Reading}
The following reputable sources provide further information for healthcare students and patients seeking reliable, current advice about food groups and a balanced diet.
\begin{itemize}
    \item National Health and Medical Research Council. \textit{Australian Dietary Guidelines.} https://www.eatforhealth.gov.au/
    \item Healthdirect Australia. \textit{Healthy eating.} https://www.healthdirect.gov.au/healthy-eating
    \item Nutrition Australia. \textit{Healthy eating pyramid and resources.} https://nutritionaustralia.org/
    \item Dietitians Australia. \textit{Food and nutrition advice.} https://dietitiansaustralia.org.au/
    \item NHS. \textit{The Eatwell Guide.} https://www.nhs.uk/live-well/eat-well/
    \item World Health Organization. \textit{Healthy diet.} https://www.who.int/
\end{itemize}
